% Generated draft: requires mathematical and source review.
\documentclass[11pt]{article}
\usepackage[margin=25mm]{geometry}
\usepackage{fontspec}
\setmainfont{Latin Modern Roman}
\usepackage{amsmath,amssymb,mathtools,unicode-math}
\setmathfont{Latin Modern Math}
\usepackage{enumitem,needspace,xurl,hyperref}
\setlength{\parindent}{0pt}
\setlength{\parskip}{5pt}
\setlength{\emergencystretch}{3em}
\newcommand{\N}{\mathbb{N}}
\newcommand{\Z}{\mathbb{Z}}
\newcommand{\Q}{\mathbb{Q}}
\newcommand{\R}{\mathbb{R}}
\newcommand{\C}{\mathbb{C}}
\newcommand{\F}{\mathbb{F}}
\newcommand{\A}{\mathbb{A}}
\newcommand{\PP}{\mathbb{P}}
\newcommand{\E}{\mathbb{E}}
\newcommand{\eps}{\varepsilon}
\newcommand{\dd}{\,\mathrm{d}}
\begin{document}
\section*{UnsolvedMath ID: 30004638}
\textbf{Source problem code:} catalogue entry
\subsection{Definitions and mathematical statement}
The supplied record denotes a family of network optimal-transport quantities by \(W_\kappa^2\) and refers to associated limiting Wasserstein and Fisher--Rao metrics. It does not define the network, the space of transported states, the transport action or distance, the parameter \(\kappa\), or the limiting metrics. In particular, the notation alone does not determine whether \(W_\kappa^2\) denotes a squared distance or the quantity intended to define the gradient-flow geometry.

The stated research objective is to analyze gradient flows associated with this family of network transport geometries and with the Wasserstein and Fisher--Rao geometries described as its limits.

This objective remains under-specified: the record supplies no energy functional, admissible initial states, boundary or vertex conditions, notion of gradient flow, parameter range, or limiting regime. It also does not say whether the requested analysis concerns existence, uniqueness, evolution equations, qualitative behavior, or convergence of gradient flows as the geometry changes. Consequently, the supplied evidence does not determine a fully quantified mathematical assertion. These missing choices must be recovered from the cited report before a precise theorem or conjecture can be formulated.
\subsection{Short English statement}
Study gradient flows for the network transport quantities denoted by \(W_\kappa^2\) and for their stated Wasserstein and Fisher--Rao limits. The supplied record leaves the geometries, energies, and intended conclusions unspecified, so this is a research objective requiring further source clarification.
\subsection{Sources}
\begin{itemize}
\item Local UnsolvedMath registry, permanent ID \(30004638\), title ``Gradient Flows for Network Transport Metrics,'' source problem code \texttt{catalogue entry}, collection ``Oberwolfach Reports Open Problems.'' The supplied proposed year is 2021; no individual proposer is identified.
\item Supplied source citation: \emph{Applications of Optimal Transportation in the Natural Sciences} (2022), Oberwolfach Reports, pp.~507--587, DOI \texttt{10.4171/owr/2021/10}. The literature-source label is ``Oberwolfach Report 10/2021.'' \url{https://doi.org/10.4171/owr/2021/10}.
\item Supplied literature reference: ``Metric gradient flows overview.'' \url{https://doi.org/10.1007/s13373-017-0101-1}. No authors, article title beyond this label, or relevant result are supplied.
\item Literature assessment checked at \texttt{2026-08-22}: the supplied curation labels the problem open and states that the report leaves gradient flows for the resulting graph transport and Fisher--Rao type metrics as an interesting open problem, with related metric theories not settling the general question. This is supplied dated curation, not an independently verified assessment; the record expressly requests independent verification of time-sensitive status and cited claims.
\end{itemize}
\end{document}
